\section{Guest-lecture Principles}

\subsection{Incremental Refinement of Research Goals}

Although research differs from software engineering, both begin with considerable ambiguity. Requirements engineering recognizes that requirements are often initially vague, or ambiguous and should therefore be elicited and refined incrementally rather than defined all at once. In software engineering, this ambiguity mainly concerns what stakeholders need from a system. In research, it concerns how a broad and underexplored problem should be defined.

My research topic is understanding football through AI, which can also be interpreted in many ways. The focus could be an individual player, a groups of players, or the complete scene on the pitch. ``Understanding'' could refer to tactical patterns, player actions, physical condition, capability, or other aspects of performance. However, these possibilities are not equally feasible or valuable. Because my research is application-oriented, the needs of the football community should also influence how the goal is refined. Relevant outcomes may provide value to coaches, players, clubs, fans, or data companies. This resembles requirements engineering, in which stakeholders and their goals are identified before specific system requirements are developed.

Research differs because the feasibility and value of a goal cannot always be known in advance. An investigation may reveal that a question is too broad, that the available data are insufficient, or that another direction is more promising. Incremental refinement is therefore particularly useful: I can begin with questions that are achievable with the available data, and then adjust the research goal as new knowledge is gained. As the project progresses, the results make the next meaningful research step clearer.

\subsection{Agentic AI in AI Research}


Agentic AI is transforming software engineering by automating tasks at multiple levels, from implementation and testing to workflow management and decision support. This automation is changing the role of software engineers, shifting their focus from routine execution to higher-level responsibilities such as defining requirements, system design, and strategic planning.
A similar evolution is occurring in AI research. From my experience, AI agents like Codex are already capable of implementing well-defined ideas, managing experimental pipelines, and summarizing results. By handling much of the repetitive engineering work, these tools enable researchers to test and iterate on ideas with greater efficiency.

However, the capabilities of current AI agents are limited when it comes to the core aspects of scientific discovery. They still lack the judgment required to determine the novelty of a research direction, identify the most insightful next step, or assess the elegance and true value of an idea. Therefore, researchers must currently define the direction, formulate clear ideas, and evaluate the scientific significance, while AI supports implementation and experimentation.
While these technological limitations may change over time, I believe the role of human direction will remain fundamental. The ultimate goal of AI research is to create tools that serve human needs and values. Therefore, it is essential that humans continue to define which problems are worth solving and which outcomes are truly meaningful.


